% =====================================================================
\section{Dữ liệu, cách chia tập và chống rò rỉ}
\label{sec:data}

\subsection{Visual Genome $\cap$ COCO (cấp độ 1)}
\label{sec:data-vg}

Tập công khai thứ nhất là phần giao giữa Visual Genome~\cite{krishna2017vg} (VG, ảnh có scene graph gán tay: vật thể,
thuộc tính, quan hệ) và MS-COCO~\cite{lin2014coco} (mỗi ảnh có 5 caption). Hai bộ nối với nhau theo \texttt{coco\_id}. Truy
vấn là caption COCO; đáp án là ảnh chứa caption đó. Nhóm chia theo split Karpathy~\cite{karpathy2015deep}, cách chia chuẩn
của bài toán truy vấn ảnh--văn bản trên COCO, và gộp phần \texttt{restval} vào train. Nhóm lấy toàn bộ phần giao, không
lấy mẫu, và bỏ 290 ảnh VG trùng \texttt{coco\_id} với ảnh khác (Bảng~\ref{tab:data}). Khi đánh giá một split, tập ứng viên
(pool) là toàn bộ ảnh của chính split đó. Pool test (2\,138 ảnh) nhỏ hơn MSCOCO 5K, nên số của nhóm không so trực tiếp được
với số trong các bài báo.

\textbf{Đồ thị ảnh.} Mỗi ảnh có một scene graph riêng, lấy từ \texttt{objects.json} và \texttt{relationships.json} của VG.
Nhãn giữ dạng mở, chỉ chuẩn hoá alias, không lọc về VG150 (bước kiểm tra độ phủ cho thấy lọc VG150 làm mất các nhãn hay
gặp như \emph{water}, \emph{grass}, \emph{field}). Quan hệ có hướng; các quan hệ trùng cả ba thành phần được gộp. Không có
cạnh giữa các ảnh, nên đồ thị không chứa cạnh suy ra từ thống kê toàn tập. Mỗi ảnh train và val có trung bình 26,7 node và
18,0 cạnh; 3,0\% ảnh không có cạnh nào.

\textbf{Đồ thị câu.} Mô hình cấp độ~1 tách caption bằng SceneGraphParser~\cite{schuster2015sgparser} (gói
\texttt{sng\_parser}): trung bình 3,06 vật thể và 1,65 quan hệ mỗi câu, và 14,2\% câu val không tách được quan hệ nào. Mô
hình chính dùng FACTUAL-T5 large~\cite{li2023factual}, một parser seq2seq đã huấn luyện sẵn để sinh scene graph từ câu.
Parser chỉ chạy suy luận, không tinh chỉnh. Trên val, FACTUAL giảm tỉ lệ node là đại từ từ 3,9\% xuống 0,1\%, và tăng tỉ lệ
câu có ít nhất một quan hệ từ 85,8\% lên 97,3\%. (Nguồn: \file{experiments/data_stats.md},
\file{experiments/improve/parser-compare/REPORT.md}.)

\begin{table}[t]
\centering
\caption{Kích thước hai tập công khai. Mỗi thực thể MKG-W có đúng một ảnh và một truy vấn.}
\label{tab:data}
\small
\begin{tabular}{lrrr}
\toprule
 & Train & Val & Test \\
\midrule
VG $\cap$ COCO: ảnh & 46\,944 & 2\,126 & 2\,138 \\
VG $\cap$ COCO: caption (truy vấn) & 234\,854 & 10\,633 & 10\,696 \\
\midrule
MKG-W retrieval: thực thể & 7\,136 & 892 & 892 \\
MKG-W KGC: triple & 34\,196 & 4\,276 & 4\,274 \\
\bottomrule
\end{tabular}
\end{table}

\subsection{MKG-W (cấp độ 2)}
\label{sec:data-mkgw}

Tập công khai thứ hai là MKG-W~\cite{xu2022mmrns}, một đồ thị tri thức đa phương thức lấy từ Wikidata~\cite{vrandecic2014wikidata}:
mỗi thực thể có nhãn, mô tả và ảnh. Nhóm dùng nó cho hai việc trên cùng một bộ thực thể.
\begin{itemize}
\item \textbf{Đồ thị tri thức.} Bản MKG-W của NativE có 15\,000 thực thể, 169 quan hệ, và các triple chia sẵn cho bài toán
link prediction (KGC) như ở Bảng~\ref{tab:data}.
\item \textbf{Bài toán truy vấn.} Chỉ xét thực thể có cả mô tả lẫn ảnh. Truy vấn là mô tả tiếng Anh của thực thể, đã che
tên, ví dụ ``a commune in Pyrénées-Atlantiques, France''. Đáp án là ảnh của chính thực thể đó. Ảnh lấy từ bản phát hành của
nhóm tác giả MMRNS: 37\,192 ảnh của 8\,961 thực thể; mỗi thực thể dùng một ảnh. Đây cũng là bộ ảnh mà NativE dùng để trích
đặc trưng ảnh. Split retrieval cố định với seed 0; tổng 8\,920 thực thể.
\end{itemize}
Trên val, 78,81\% truy vấn liên kết được với ít nhất một thực thể của đồ thị (test 80,94\%). Với 45,6\% truy vấn val, đáp án
nối trực tiếp với một thực thể được nhắc trong câu bằng một cạnh KGC-train.

\subsection{Chống rò rỉ dữ liệu}
\label{sec:leakage}

Mục~6 của đề là tiêu chí chấm. Bảng~\ref{tab:leakage} liệt kê từng quy tắc và cách nhóm bảo đảm nó trong code hoặc quy
trình.

\begin{table*}[t]
\centering
\caption{Các quy tắc chống rò rỉ và cách bảo đảm.}
\label{tab:leakage}
\small
\begin{tabular}{p{0.30\textwidth}p{0.64\textwidth}}
\toprule
Quy tắc & Cách bảo đảm \\
\midrule
Mọi mô hình so sánh dùng cùng một tệp chia &
File split trong \file{data/splits/} là nguồn duy nhất; SHA-256 ghi ở \file{data/MANIFEST.md}. Mọi dòng trong một bảng dùng
chung danh sách top-50 của CLIP. \\
Test không dùng để huấn luyện, chọn siêu tham số hay chọn mô hình &
VG: mỗi mô hình đã chốt chấm test một lần, 25/09 (mô hình cấp~1) và 25/10 (mô hình chính); scene ranker
(Mục~\ref{sec:contrib}) cũng chấm một lần. \texttt{GraphStore} từ chối nạp test nếu không gọi rõ
\texttt{allow\_test=True}; trước khi chấm, script kiểm tra ứng viên test tái tạo đúng số của CLIP thuần.
MKG-W: lần chấm báo cáo là 17/11. Trước đó, một lần chạy thử code với embedding giả trên test đã in ra số của dòng đếm
cạnh; ngày 18/11, một thành viên chạy lại độc lập toàn bộ cấp~2 với NativE tự huấn luyện (Mục~\ref{sec:l2}). Cả hai lần
không dẫn tới thay đổi nào: $\alpha,\beta$ chỉ lấy từ val. \\
Cạnh thống kê chỉ xây từ train &
Đồ thị ảnh VG là chú thích của từng ảnh, không có cạnh giữa các ảnh. Kênh đồ thị của MKG-W chỉ nạp triple KGC-train; hàm
dựng đồ thị có \texttt{assert kgc\_split == "train"}. \\
$\alpha$ và mọi lựa chọn dựa trên train/val &
VG: trọng số kết hợp $a,b,c$ học trên train; dừng sớm, chọn checkpoint, ablation và từng bước xây mô hình chính dựa trên
val; lr và weight decay cố định từ đầu, không dò. Với MKG-W, $\alpha,\beta$ chọn trên val
cho từng seed, ghi vào \file{experiments/level2/val/choice.json}; lần chạy test chỉ đọc lại file này. \\
Ứng viên lúc huấn luyện không chạm val/test &
Train được chia cố định (seed 0) thành 22 nhóm khoảng 2\,134 ảnh. Top-50 của mỗi caption train chỉ lấy trong nhóm của ảnh
đúng, nên độ khó giống lúc đánh giá. CLIP và FACTUAL-T5 đóng băng. \\
Mô hình đồ thị tri thức không thấy truy vấn và ảnh của val/test &
Đặc trưng văn bản của NativE mã hoá chính mô tả thực thể, tức chính câu truy vấn; đặc trưng ảnh mã hoá chính ảnh ứng viên.
Bản \texttt{MKG-W-rmask} đặt về 0 hai đặc trưng này cho 1\,784 thực thể val/test retrieval trước khi huấn luyện NativE. \\
Bộ truy vấn tự viết chỉ để đánh giá &
Bộ 25 câu MKG-W được commit một lần (16/11) cùng script kiểm tra, không sửa sau đó; lần chấm dùng lại $\alpha,\beta$
của val. \\
\bottomrule
\end{tabular}
\end{table*}

\subsection{Kiểm tra pipeline}
\label{sec:checks}

Trước khi tin các con số, nhóm kiểm tra pipeline bằng những phép so có đáp án biết trước.
\begin{itemize}
\item \textbf{CLIP trên MSCOCO 5K.} Chạy chính code của nhóm (mã hoá, FAISS, module đánh giá) trên pool chuẩn 5\,000 ảnh.
Kết quả lệch số đã công bố của open\_clip~\cite{cherti2023openclip} không quá 0,03 điểm (text$\to$ảnh R@1 30,45 so với
30,44). Phép so này cũng phát hiện một lỗi: trọng số OpenAI phải nạp bằng cấu hình \texttt{ViT-B-32-quickgelu}. Nạp sai
cấu hình làm R@1 thấp hơn khoảng 1 điểm. Toàn bộ đặc trưng được trích lại sau khi sửa
(\file{experiments/checks/README.md}).
\item \textbf{Mô hình chưa huấn luyện.} Ở bước khởi tạo, mô hình cấp độ~1 phải cho đúng điểm của phép thử không huấn luyện
đã đo trước (R@1 val 43,50). Có script kiểm tra tự động cho việc này.
\item \textbf{Demo.} Gõ 200 caption val vào giao diện. Demo cho cùng 50 ứng viên, cùng đồ thị câu và cùng top-1 với pipeline
đánh giá; điểm cuối lệch tối đa $4\cdot10^{-6}$.
\item \textbf{Đặc trưng MKG-W dựng lại.} Đặc trưng CLIP gốc của MKG-W không còn truy cập được, nên nhóm trích lại bằng đúng
script trong repo. Trên val, CLIP thuần cho lại đúng R@1 24,33 như kết quả đã commit trước đó.
\end{itemize}
